\documentclass{article}
\usepackage[T1]{fontenc}
\usepackage{lmodern}
\usepackage{graphicx}
\usepackage{amsmath,amssymb}
\usepackage[a4paper,margin=2cm]{geometry}
\usepackage[absolute,overlay]{textpos}
\setlength{\TPHorizModule}{1mm}
\setlength{\TPVertModule}{1mm}
\usepackage[indonesian]{babel}
\usepackage{hyperref}
\setlength{\emergencystretch}{2em}
% unit: Halaman 8

\begin{document}

% === Halaman 8 (python/native) ===

• Jika a adalah akar primitif dari bilangan prima p, maka untuk bilangan bulat b kita 

dapat menemukan pangkat x sedemikian sehingga

b ax (mod p) ,

0 x ( p – 1)

• Pangkat x disebut logaritma diskrit dari b untuk basis a (mod p).

• Dalam persoalan logaritma diskrit, diberikan b  ax (mod p), carilah x yang 

memenuhi kekongruenan tersebut.  

• Sebagai contoh, 7 adalah akar primitif dari bilangan prima  p = 41. Maka,

   carilah x sedemikian sehingga  15  7x (mod 41), jawabannya adalah 3 karena 


73 = 343  15 (mod 41) 

8

\end{document}
